\section*{Data availability}

The data that support the findings of this study are publicly available. The streamflow and meteorological data are derived from the CAMELS-US dataset\cite{CAMELS}. The dataset is accessible via: \href{https://doi.org/10.1016/0022-1694(89)90184-4}{https://doi.org/10.1016/0022-1694(89)90184-4}. 

\section*{Code availability}

The source code and processed datasets are available in the GitHub repository: \href{https://github.com/Polyethylenekjc/Sparse-Transformer-In-Streamflow-Prediction}{https://git\\hub.com/Polyethylenekjc/Sparse-Transformer-In-Streamflow-Prediction}.

\section*{CRediT authorship contribution statement}

\textbf{Jiachen Kong}: Conceptualization, Methodology, Software, Writing -- original draft.  
\textbf{Zhaocai Wang}: Methodology, Software, Validation, Supervision, Writing -- review \& editing.  
\textbf{Jun Wang}: Validation, Supervision, Writing -- review \& editing.  
\textbf{Changhyun Jun}: Supervision, Writing -- review \& editing.  
\textbf{Zhaoyang Zhu}: Visualization, Writing -- original draft.
\textbf{Chunyong Wei}: Writing - review \& editing.

\section*{Declaration of Competing Interest}

The authors declare that they have no known competing financial interests or personal relationships that could have appeared to influence the work reported in this paper.

\section*{Acknowledgments}
This study was supported by the Soft Science Project of Shanghai, China (No. 25692108900), the Humanities and Social Science Fund of Ministry of Education of China, China (No. 24YJAZH167), and the National Statistical Science Research Project of China, China (No. 2025LY082).

\begin{table}[htbp]
\centering
\caption{Appendix 1: Main Terminology Overview}
\label{tab:terminology_overview}
\begin{tabularx}{\textwidth}{p{3cm} p{2cm} X}
\toprule
\textbf{Full Name} & \textbf{Abbreviation} & \textbf{Description} \\
\midrule
Physics-Informed Neural Networks & PINNs & Optimization constraint, emphasizes that predictions must obey physical laws, e.g., streamflow predictions satisfy mass conservation. \\
Sparse Transformer for Explainable Hydrology & STEH & Architecture design, emphasizes sparse and analyzable internal paths, e.g., identifying paths responsible for evapotranspiration. \\
Circuit & Circuit & Functional subset of computational units selected by sparsification (e.g., specific heads/neurons) that jointly supports prediction. \\
Structural & Structural & Organization pattern of components and interactions within a circuit (e.g., layer-wise arrangement and dependency layout). \\
Information Flow Pathway & Pathway & Directed route of information propagation within a circuit (e.g., precipitation-to-runoff transfer through key nodes). \\
Explainable AI & XAI & Overall goal, broad transparency and interpretability, e.g., understanding why the model predicts more accurately during drought periods. \\
\bottomrule
\end{tabularx}
\end{table}

\begin{table}[htbp]
\centering
\noindent\textbf{Appendix 2: Confusable Concepts and Verification Methods}
\vspace{0.5em}
\begin{tabularx}{\textwidth}{p{4cm} X}
\toprule
\textbf{Full Name} & \textbf{Verification / Implementation Method} \\
\midrule
Explainability & Post-hoc analysis (SHAP / LIME) or intrinsic structural analysis to understand model prediction behavior. \\
Physical Structure Explainability & Structural sparsification + linear probing to verify whether internal computational paths correspond to physical quantities (e.g., soil moisture). \\
Causal Contribution & Rigorous causal verification, including causal graphs / SEMs, intervention experiments (Do-Calculus), and counterfactual / A/B comparisons to ensure true causal effect from input to output. \\
Structural Causal Contribution & Structured intervention + causal pathway verification + counterfactual comparison, to verify that internal components (attention heads, paths) are necessary and irreplaceable in the input–output causal chain. \\
\bottomrule
\end{tabularx}
\end{table}
\FloatBarrier